\section{Результаты I: способность решать задачи}
\label{sec:capability}

% Figure 2: the four-panel characterisation of the method.
% Every coordinate is read from figures/data/*.dat, which analysis/build_paper.py
% regenerates from the canonical run CSV. No number is typed here.

\begin{figure*}[t]
\centering
\begin{tikzpicture}[font=\small]

\begin{axis}[
  name=A,
  width=0.46\textwidth, height=4.6cm,
  ybar, bar width=15pt,
  ymin=0, ymax=0.75,
  ylabel={доля решённых},
  xtick={0,1,2,3},
  xticklabels={контроль,$k{=}0$,$k{=}2$,$k{=}6$},
  title={\textbf{(a)} способность решать задачи},
  enlarge x limits=0.15,
  grid=major, grid style={black!12},
  error bars/y dir=both, error bars/y explicit,
]
\addplot[draw=black!60, fill=blue!25] table[x=x, y=rate,
  y error plus=ehi, y error minus=elo] {figures/data/solve_rate.dat};
\end{axis}

\begin{axis}[
  name=B, at={($(A.east)+(1.7cm,0)$)}, anchor=west,
  width=0.46\textwidth, height=4.6cm,
  ybar, bar width=9pt,
  ymin=0, ymax=95,
  ylabel={несовпадающие задачи},
  xtick={0,1,2}, xticklabels={$k{=}0$,$k{=}2$,$k{=}6$},
  title={\textbf{(b)} парное расхождение},
  enlarge x limits=0.25,
  grid=major, grid style={black!12},
  legend style={at={(0.97,0.97)}, anchor=north east, font=\scriptsize,
                draw=black!25, fill=white, fill opacity=0.9, text opacity=1},
]
\addplot[draw=black!60, fill=red!35] table[x=x, y=ctlonly]
  {figures/data/discordance.dat};
\addlegendentry{решила только контрольная группа}
\addplot[draw=black!60, fill=green!35] table[x=x, y=armonly]
  {figures/data/discordance.dat};
\addlegendentry{решила только рука состояния}
\end{axis}

\begin{axis}[
  name=C, at={($(A.south)-(0,1.85cm)$)}, anchor=north,
  width=0.46\textwidth, height=4.6cm,
  ybar, bar width=9pt,
  ymin=0, ymax=52,
  ylabel={чистые потерянные задачи},
  xtick={0,1,2}, xticklabels={$k{=}0$,$k{=}2$,$k{=}6$},
  title={\textbf{(c)} потеря по горизонту},
  enlarge x limits=0.25,
  grid=major, grid style={black!12},
  legend style={at={(0.97,0.97)}, anchor=north east, font=\scriptsize,
                draw=black!25, fill=white, fill opacity=0.9, text opacity=1},
]
\addplot[draw=black!60, fill=orange!30] table[x=x, y=short]
  {figures/data/horizon_deficit.dat};
\addlegendentry{короткие задачи}
\addplot[draw=black!60, fill=purple!45] table[x=x, y=long]
  {figures/data/horizon_deficit.dat};
\addlegendentry{длинные задачи}
\end{axis}

\begin{axis}[
  name=D, at={($(B.south)-(0,1.85cm)$)}, anchor=north,
  width=0.46\textwidth, height=4.6cm,
  xmin=0.35, xmax=1.12, ymin=0.85, ymax=2.85,
  xlabel={пиковый промпт, отношение к контролю},
  ylabel={стоимость, отношение к контролю},
  title={\textbf{(d)} чего стоит экономия промпта},
  grid=major, grid style={black!12},
  nodes near coords={\scriptsize\pgfplotspointmeta},
  point meta=explicit symbolic,
  every node near coord/.append style={anchor=west, xshift=3pt},
]
\addplot[only marks, mark=*, mark size=2.6pt, draw=black!70, fill=blue!45]
  table[x=peak, y=cost, meta=label] {figures/data/prompt_vs_cost.dat};
\draw[dashed, black!45] (axis cs:0.35,1) -- (axis cs:1.12,1);
\node[font=\scriptsize, text=black!55, anchor=south east]
  at (axis cs:1.10,1.02) {паритет};
\end{axis}

\end{tikzpicture}
\caption{Характеристика памяти со структурированным состоянием на \EOneTasks{}
программных задачах, \EOneRuns{} запусков. \textbf{(a)} Доля решённых задач
с интервалами Уилсона; контрольная группа решает \SolveBase{} против
\SolveTailZero{} для опубликованной конфигурации. \textbf{(b)} Парное
расхождение против контрольной группы того же испытания — только
несовпадающие задачи несут информацию, и они подавляюще односторонни при
$k{=}0$. \textbf{(c)} Чистые потерянные задачи, разбитые по медианному
горизонту контрольной группы: дефицит \emph{расширяется} на длинных
задачах, инвертируя аргумент в пользу метода. \textbf{(d)} При равном исходе
опубликованная конфигурация вдвое сокращает пиковый промпт и умножает
стоимость на \CostRatioTailZero. Восстановление шести обменов ($k{=}6$)
устраняет дефицит способности решать задачи и сокращает премию по стоимости
с \CostRatioTailZero$\times$ до \CostRatioTailSix$\times$, но не опускается
ниже контрольной группы: при равном исходе она остаётся более дорогой рукой
($p={}$\CostSignPTailSix).}
\label{fig:fourpanel}
\end{figure*}


Рисунок~\ref{fig:fourpanel} обобщает этот и два следующих раздела; таблицы
ниже приводят те же величины в точности.

\subsection{Доля решённых задач}

Таблица~\ref{tab:solve} даёт непарную картину. Опубликованная конфигурация
решает \SolveTailZero{} запусков против \SolveBase{} у контроля, и
доверительные интервалы не пересекаются. Возвращение истории поднимает долю
постепенно: \SolveTailTwo{} при двух обменах, \SolveTailSix{} при шести.

\begin{table}[t]
\centering
\caption{Доля решённых задач по условиям. \EOneRuns{} запусков,
\EOneTasks{} задач $\times$ \EOneTrials{} испытаний. Интервалы — по Уилсону.}
\label{tab:solve}
\input{tables/t1_solve_rate}
\end{table}

\subsection{Парное сравнение}

Непарные доли недооценивают то, что может этот дизайн: каждое условие
прошло каждую задачу, поэтому информативно не само значение доли, а
множество задач, на которых две руки расходятся. Таблица~\ref{tab:mcnemar}
сравнивает каждую руку с контролем того же испытания.

\begin{table}[t]
\centering
\caption{Парная способность решать задачи против контроля того же испытания,
точный тест Макнемара. Информацию несут только несовпадающие ячейки.}
\label{tab:mcnemar}
\input{tables/t2_mcnemar}
\end{table}

Опубликованная конфигурация проигрывает \McnCtlTailZero{} задач и выигрывает
\McnArmTailZero{} ($p={}$\McnPTailZero). Эффект велик, односторонен и далеко
за порогом значимости. При двух сохранённых обменах дефицит всё ещё ясен
(\McnCtlTailTwo{} против \McnArmTailTwo, $p={}$\McnPTailTwo); при шести он
исчезает: \McnCtlTailSix{} против \McnArmTailSix{} неотличимо от шума
($p={}$\McnPTailSix).

Эта последняя строка — самая информативная в статье. Шесть дословных обменов
почти ничего не добавляют к промпту, но возвращают способность решать задачи
к паритету. Значит, ущерб наносят не суммирование состояния и не его
структурирование, а отбрасывание недавнего транскрипта. То, что блок
состояния отнимает у решаемости, возвращает маленький кусок истории рядом с
ним, а не сам блок.

\subsection{Горизонт: потеря растёт там, где была обещана выгода}

Задачи разбиты по медианному числу ходов контроля (\HorizonCut{} ходов):
получается \HorizonShort{} коротких и \HorizonLong{} длинных задач.
Горизонт берётся из \emph{контрольной} руки: считать по числу ходов самой
экспериментальной руки было бы циклично, потому что рука, которая заново
добывает данные, по построению делает больше ходов.

\begin{table}[t]
\centering
\caption{Парная способность решать задачи по стратам горизонта, разбитым по
медиане контроля в \HorizonCut{} ходов.}
\label{tab:horizon}
\input{tables/t4_horizon}
\end{table}

У опубликованной конфигурации дефицит на коротких задачах составляет
\HorTailZeroShortCtl{} против \HorTailZeroShortArm{}
($p={}$\HorTailZeroShortP), а на длинных вырастает до
\HorTailZeroLongCtl{} против \HorTailZeroLongArm{} ($p={}$\HorTailZeroLongP).

Это переворачивает аргумент в пользу метода. Структурированное состояние
мотивируют длинными горизонтами, на которых неограниченный транскрипт должен
якобы стать неуправляемым. На этом корпусе длинный горизонт — как раз то
место, где отбрасывание транскрипта обходится дороже всего: дефицит
накапливается по ходу задачи в ту же сторону, в какую ожидали накопления
выгоды.

При $\tailk{=}6$ ни одна страта различия не показывает, поэтому эффект
горизонта — свойство того, сколько истории было удалено, а не свойство блока
состояния.

\subsection{Как завершались запуски}

Сравнение решаемости честно, только если ограничение харнеса не
засчитывается агенту как неспособность. Таблица~\ref{tab:stop} разделяет
четыре терминальных состояния.

\begin{table}[t]
\centering
\caption{Причины остановки по условиям. Остановки по бюджету
сосредоточены в руке, которая тратит больше всего ходов, поэтому они
исключены из агрегации стоимости.}
\label{tab:stop}
\input{tables/t7_stop_reasons}
\end{table}

Исчерпание бюджета ходов встречается в руках с состоянием и никогда — в
контроле и при $\tailk{=}6$\live. Это само по себе свидетельство в пользу
механизма из \S\ref{sec:mechanism}: рука, вынужденная заново добывать свой
контекст, — это та же рука, у которой кончаются ходы. Это же и заранее
объявленная ловушка из \S\ref{sec:design}, и потому каждое утверждение о
стоимости в \S\ref{sec:economics} вычисляется только по парам, где задачу
решили обе руки.
